	\documentclass[12pt,a4paper,italian]{article}


\usepackage[italian]{babel}
\usepackage[latin1]{inputenc}
\usepackage{amsmath}
\usepackage{amsfonts}
\usepackage{amssymb}
\usepackage{color}
\usepackage{xcolor}
\usepackage{hyperref}
\usepackage[all]{hypcap}
\usepackage{ifthen}
\usepackage{wrapfig}
%\author{Piero Bizzotto}
\usepackage[top=2cm,bottom=5cm,left=80pt,right=80pt]{geometry}
\usepackage{graphicx}
\DeclareGraphicsExtensions{.jpg,.png}

\newcommand{\ajax}{AJAXDRAW}
\newcommand{\sito}{\href{http://ajaxdraw.sourceforge.net}{http://ajaxdraw.sourceforge.net}}

\setlength{\parindent}{0pt} %settato indentazione di default 
\setlength{\headheight}{3cm} %settato grandezza header...in altre parole, quanto distanzio il doc dall'intestazione

\usepackage{fancyhdr} %pacchetto per le intestazioni
\pagestyle{fancy} %uso del pacchetto


\fancyhead{} %annulla head di default
\fancyfoot{} %annulla foot di default


\usepackage{lastpage} %setto pg di pgtot a rfoot
     \rfoot{pagina \thepage\ di \pageref{LastPage}}


\lfoot{Versione: \insertversion} %setto versione doc a lfoot
\renewcommand{\footrulewidth}{0.5pt} %ridefinisco il valore della riga di intestazione
\renewcommand{\headrulewidth}{0.5pt} %ridefinisco il valore della riga di pie' di pagina

\newcommand{\insertversion}{0.0} %definisco il nuovo comando per inserire la versione


\lhead{  \begin{Huge} \ajax \end{Huge} \\  %intestazione di sinistra
					%\begin{Large}	Software per il Disegno Grafico\\ in Tecnologie Web \end{Large}  
					\begin{normalsize}\sito \end{normalsize}
			%\\ versione documento: \insertversion\ del \today} %setto l'intestazione sx
		}
\rhead{ %includo logo nell'intestazione dx
	 	\includegraphics[scale=0.5]{../logo/logo}  
}


%CREAZIONE ELENCHI NUMERATI PERSONALIZZATI
\newcounter{Lcount}
\newcounter{Rcount}
\setcounter{Lcount}{0}
\setcounter{Rcount}{0}

\newenvironment{elenconumerato}[2][ ]
{
  \begin{list}{#1\arabic{Lcount}.}
    {
	\setcounter{Rcount}{\value{Lcount}}
	\setcounter{Lcount}{0} 
	\usecounter{Lcount} 
\addtolength{\leftmargin}{#2pt}
	}
}
{
  \end{list}
 \setcounter{Lcount}{\value{Rcount}}
}

%CREAZIONE ELENCHI PUNTATI
\newenvironment{elencopuntato}[1][]
{
\begin{list}{\textbullet} %\itemindent=#1pt
	{
	\addtolength{\leftmargin}{#1pt}
	}
} 
{
\end{list}
}


\newenvironment{elencodescrittivo}[1][]{\begin{description} \setlength{\itemindent}{#1pt} \addtolength{\leftmargin}{#1pt}} {\end{description}}

\newcommand{\TITOLODOC}{Titolo}

%footer centrale
\cfoot{ \TITOLODOC \\  E-mail:    \href{ mailto:webshape.contact@gmail.com}{ webshape.contact@gmail.com}  }

%INSERIMENTO IMMAGINI
\newcommand{\imagerealsize}[1]{\vspace{20pt} \includegraphics{#1} }
\newcommand{\imageadapted}[1]{\vspace{20pt} \includegraphics[width=1\textwidth]{#1} }

\newcommand{\glosspath}{.\glossario}
\newcommand{\gloss}[1]{\hyperref{\glosspath~\glossario.pdf}{}{#1}{#1}}

\hypersetup{
    %bookmarks=true,         % show bookmarks bar?
    %unicode=false,          % non-Latin characters in Acrobat’s bookmarks
	%pdftoolbar=true,        % show Acrobat’s toolbar?
	%pdfmenubar=true,        % show Acrobat’s menu?
    %pdffitwindow=true,      % page fit to window when opened
    %pdftitle={My title},    % title
    %pdfauthor={Author},     % author
    %pdfsubject={Subject},   % subject of the document
    %pdfnewwindow=true,      % links in new window
    %pdfkeywords={keywords}, % list of keywords
    colorlinks=true,         % false: boxed links; true: colored links
    linkcolor=black,           % color of internal links
    %citecolor=green,        % color of links to bibliography
    %filecolor=magenta,      % color of file links
    urlcolor=teal    % color of external links
%	linktocpage=false;
}


%COLORAZIONE TESTO
\newcommand{\blue}[1]{{\color {blue} #1}} 
\newcommand{\red}[1]{{\color {red} #1}}
\newcommand{\green}[1]{{\color {green} #1}}
\newcommand{\sezione}[1]{\leftskip=0pt \section{#1} \leftskip=18pt}
\newcommand{\subsezione}[1]{\leftskip=18pt \subsection{#1} \leftskip=36pt}
\newcommand{\subsubsezione}[1]{\leftskip=36pt \subsubsection{#1} \leftskip=54pt}
\newcommand{\subsubsecindent}{54}
\newcommand{\subsecindent}{36}
\newcommand{\secindent}{18}
\newcommand{\normindent}{8}
\newcommand{\code}[1]{{\bfseries \texttt{#1}}}
\newcommand{\paragrafo}[1]{\leftskip=36pt \paragraph{#1} \leftskip=54pt}
\newcommand{\subparagrafo}[1]{\leftskip=54pt \subparagraph{#1} \leftskip=72pt} %BASE!!!

\title{\TITOLODOC}
\author{Dissegna Stefano}

\begin{document}

\renewcommand{\insertversion}{1.0} %INSERIRE LA VERSIONE QUI DENTRO STILE x.x.xx
\renewcommand{\TITOLODOC}{Piano di qualifica} %INSERIRE IL TITOLO DEL DOCUMENTO DA FAR COMPARIRE A PIE PAGINA
\renewcommand{\glosspath}{.\glossario} %INSERIRE PERCORSO RELATIVO

%%%%%%%%%%%%%%%%%%%%%%PARTE DA NON MODIFICARE%%%%%%%%%%%%%%%%%
\begin{titlepage}
\begin{center}
	\begin{Large}	\today \end{Large}
\end{center}

\vspace{20pt}

\begin{center}
	\begin{Huge}
				\textbf{\ajax}
	\end{Huge}
\end{center}			

\begin{center}
	\begin{large}
				\textbf{Software per il Disegno Grafico\\ in Tecnologie Web}
	\end{large}
\end{center}			

\vspace{20pt}

\begin{center}
\includegraphics[width=150pt]{../logo/logo}
\end{center}

\vspace{170pt}
\begin{center} %INSERIRE ALL'INTERNO IL TITOLO DOCUMENTO CHE COMPARIRA NELLA PAGINA INIZIALE				
	\begin{Huge}
				\textbf{\TITOLODOC}
	\end{Huge}
			\\
\end{center}
\vspace{210pt}
\begin{center}
Versione: \insertversion
\end{center}
\end{titlepage}

\newpage
%%%%%%%%%%%%%%%%%%%%%%FINE PARTE DA NON MODIFICARE%%%%%%%%%%%%%%%%%

\begin{center} %INSERIRE ALL'INTERNO IL TITOLO DOCUMENTO CHE COMPARIRA NELLA PAGINA INIZIALE
	\begin{Huge}	
				\textbf{\TITOLODOC}
			\\
	\end{Huge}
\end{center}

%\setlength{\parindent}{18pt} %settato indentazione di default 
\section*{\LARGE Sommario:} %SEZIONE SOMMARIO
\indent \indent
Il presente documento descrive le strategie per la gestione della qualit\`a adottate dall'azienda WebShape, e comprende un resoconto delle relative attivit\`a effettivamente svolte.

\section*{\LARGE Stato del documento:}
\indent \indent
	Formale Esterno

\section*{\LARGE Redazione:}
	\begin{table}[!h]
		\begin{center}
			\begin{tabular}
				{|c|c|}
				\hline
				%%%%%%%%%%%%%%INTESTAZIONE COLONNE%%%%%%%%%%%%%%%%%%%%%%%%%%%%%%%%
				\multicolumn{2}{|c|}{ \textbf{Redazione} } \\
				\hline
				\textbf{Fase} & \textbf{Redattori} \\
				%%%%%%%%%%%%%%FINE INTESTAZIONE COLONNE%%%%%%%%%%%%%%%%%%%%%%%%%%%%%%%%%%%%%%
				\hline
				%%%%%%%%%%% PARTE DA MODIFICARE %%%%%%%%%%%%%%%%%%%%%%%%%%%%%%%%%%%%%%%%%%		
				Pre-RR & Dissegna Stefano\\
				\hline
				RR-RPP & \\
										& \\
				\hline
				%%%%%%%%%%% FINE PARTE DA MODIFICARE %%%%%%%%%%%%%%%%%%%%%%%%%%%
			\end{tabular}
			\caption{Lista Redattori} %INSERIRE DIDASCALIA - SE NECESSARIA - 
			\label{tabredazione}
		\end{center}
	\end{table}	

\newpage
\section*{\LARGE Approvazione:}
\begin{table}[!h]
	\begin{center}
		\begin{tabular}
			{|c|c|}
			\hline
			%%%%%INTESTAZIONE COLONNE%%%%%%%%%%%%%%%%%%%%%%%%%%%%%%%
			\multicolumn{2}{|c|}{ \textbf{Approvazione} } \\
			\hline
			\textbf{Fase} & \textbf{Approvatori} \\
			%%%%%%%%%%%%%%FINE INTESTAZIONE COLONNE%%%%%%%%%%%%%%%%%%%%%%%%%%%%%%
			\hline
			%%%%%%%%%%% PARTE DA MODIFICARE %%%%%%%%%%%%%%%%%%%%%%%%%%%%%%%%%%%%%%		
			Pre-RR  &  Dal Bosco Davide\\
									&  \\
			\hline
			RR-RPP & \\
									& \\
			\hline
			%%%%%%%%%%% FINE PARTE DA MODIFICARE %%%%%%%%%%%%%%%%%%%%%%%%%%%%%%%%%%%
		\end{tabular}
		\caption{Lista Approvatori} %INSERIRE DIDASCALIA - SE NECESSARIA - 
		\label{tabapprovazione}
	\end{center}
\end{table}


\section*{\LARGE Lista di Distribuzione:}

	\begin{elenconumerato}{\normindent}
		\item WebShape
		\item I committenti Conte Renato e Vardanega Tullio in rappresentanza \\  dell'azienda proponente Zucchetti SPA
	\end{elenconumerato}

\newpage



\section*{\LARGE Registro delle Modifiche:}


\begin{center}
	\begin{table}[h]
		  \begin{tabular*}
			{1\textwidth}%
				{@{\extracolsep{\fill}}|p{0.1\textwidth}|p{0.54\textwidth}|p{0.26\textwidth}|}
			 \hline
%%%%%%%%%%%%%%INTESTAZIONE COLONNE%%%%%%%%%%%%%%%%%%%%%%%%%%%%%%%%%%%%%%%%%%%%%%%%%%%%%%%%%%%%%%
			\textbf{Versione}  & \textbf{Descrizione} & \textbf{Autore} \\
%%%%%%%%%%%%%%FINE INTESTAZIONE COLONNE%%%%%%%%%%%%%%%%%%%%%%%%%%%%%%%%%%%%%%%%%%%%%%%%%%%%%%%%%%%%%%
		 \hline
%%%%%%%%%%% PARTE DA MODIFICARE %%%%%%%%%%%%%%%%%%%%%%%%%%%%%%%%%%%%%%%%%%%%%%%%%%%%%%%%%%%%%%%%%
			1.0 & 	 09$\slash$12$\slash$2008 Verifica finale in preparazione al rilascio & Dissegna Stefano\\
			\hline
             0.3 & 7/12/2008 Correzioni in seguito a revisione & Dissegna Stefano \\
			 \hline
            0.2 & 7/12/2008 Impaginazione e modifiche minori & Dissegna Stefano \\
		 	\hline
            0.1 & 23/11/2008 Piano di qualifica per la progettazione & Dissegna Stefano \\ 
            \hline
            0.0 & 22/11/2008 Bozza iniziale & Dissegna Stefano \\
             \hline 
%%%%%%%%%%% FINE PARTE DA MODIFICARE %%%%%%%%%%%%%%%%%%%%%%%%%%%%%%%%%%%%%%%%%%%%%%%%%%%%%%%%%%%
		\end{tabular*}
	\caption{tabella delle modifiche} %INSERIRE DIDASCALIA - SE NECESSARIA - 
	\label{tab:modifiche}
	\end{table}
\end{center}


\newpage
\thispagestyle{fancy}
\tableofcontents
\thispagestyle{fancy}
\newpage

\sezione{Introduzione}

\subsezione{Scopo del documento}
Il presente documento definisce il livello di qualit\`a atteso del prodotto e del processo di realizzazione dello stesso, indica le metodologie da adottare per assicurare la qualit\`a e riporta le attivit\`a svolte inerenti la gestione della qualit\`a.

\subsezione{Scopo del prodotto}
Il prodotto intende permettere la realizzazione di grafica vettoriale tramite una comoda interfaccia web.

\subsezione{Glossario}
Si veda il \textit{Glossario}.

\subsezione{Riferimenti normativi}
Il presente documento \`e redatto in accordo con le norme interne di WebShape, raccolte nel documento \textit{NormeDiProgetto} consegnato assieme a questo documento, e consultabile inoltre dal repository pubblico al quale WebShape si appoggia per i suoi progetti.

\sezione{Visione generale della strategia di verifica}

\subsezione{Organizzazione, pianificazione strategica e temporale, responsabilit\`a}
Le attivit\`a di verifica si dovranno svolgere immediatamente in seguito ad una modifica, o ad un insieme di modifiche. Qualora sia possibile svolgere l'attivit\`a in modo automatico, essa dovr\`a essere applicata in seguito ad ogni singolo cambiamento. I test automatici saranno predisposti dai verificatori. Negli altri casi, onde evitare di attivare un'intero ciclo di verifica con gli annessi costi per delle modifiche di scarso rilievo, l'attivit\`a sar\`a svolta quando le differenze tra la versione attuale del prodotto e la versione su cui \`e stata applicata l'ultima verifica sono sufficientemente rilevanti. \`E compito di chi effettua le modifiche di segnalare ai verificatori la necessit\`a di un'attivit\`a di verifica. Il responsabile della qualifica dovr\`a accertarsi che il tutto si svolga correttamente e che segua quanto descritto nel presente documento.
\paragraph{Verifica dei documenti} Il revisore, attraverso un'attenta lettura del documento, deve controllare, in ordine crescente di importanza, che:
\begin{elenconumerato}[\textbf{}]{\subsubsecindent}
\item la grammatica e la sintassi siano corrette
\item la formattazione del documento sia corretta e segua gli standard contenuti nelle \textit{NormeDiProgetto}
\item i concetti vengano espressi in modo chiaro e non ambiguo
\item i concetti espressi siano corretti
\item il documento sia completo
\item il documento risolva il problema corretto
\end{elenconumerato}
\paragraph{Verifica del progetto architetturale}
Il progetto architetturale dovr\`a essere controllato secondo le seguenti caratteristiche:
\begin{elenconumerato}[\textbf{}]{\subsubsecindent}
\item Aderenza alle \textit{NormeDiProgetto}.
\item Completezza: per controllare l'aderenza ai requisiti, dovr\`a essere verificata la tracciabilit\`a tra ogni requisito e la parte, o le parti, che lo soddisfano nel progetto.
\item Sufficienza: ogni elemento del progetto dovr\`a, o soddissfare un requisito, o essere una dipendenza, diretta o indiretta, di un elemento necessario a soddisfare un requisito, in modo che nessuna parte del progetto risulti superflua.
\item Estendibilit\`a e manutenibilit\`a: il progetto dovr\`a prevedere la possibilit\`a che il prodotto possa essere facilmente esteso e mantenuto in futuro.
\item Efficienza: sebbene si ritenga che l'ottimizzazione prematura vada evitata il pi\`u possibile, considerata la necessit\`a da parte dell'applicazione di un server, la fase di progettazione dovr\`a assicurare che il server sia in grado di gestire un numero elevato di client contemporaneamente senza che l'utente noti problemi di performance.
\end{elenconumerato}
\paragraph{Verifica del progetto di dettaglio}
Il progetto di dettaglio dovr\`a soddisfare i seguenti aspetti:
\begin{elenconumerato}[\textbf{}]{\subsubsecindent}
\item Aderenza alle \textit{NormeDiProgetto}.
\item Completezza: ogni elemento presente nel progetto architetturale dovr\`a essere presente anche nel progetto di dettaglio. Non dovranno inoltre essere presenti delle lacune tra un elemento del progetto e la sua effettiva realizzazione.
\item Sufficienza: non dovranno essere presenti elementi non strettamente necessari alla realizzazione del progetto architetturale.
\item Applicabilit\`a: il progetto dovr\`a essere realizzabile dai programmatori usando le tecnologie previste.
\end{elenconumerato}
\paragraph{Verifica del codice}
La verifica del codice si compone dei seguenti aspetti:
\begin{elenconumerato}[\textbf{}]{\subsubsecindent}
\item Verifica automatica: verr\`a effettuata tramite test di unit\`a predisposti dai verificatori. I test dovranno coprire tutto il codice, a eccezione delle parti il cui funzionamento dipende da input esterni non simulabili in un test di unit\`a, come ad esempio il codice che gestisce gli eventi generati dal movimento del mouse, oppure da effetti secondari i cui risultati non siano facilmente confrontabili in modo automatico, ad esempio come appare l'aspetto grafico del programma sullo schermo dell'utente. I test dovranno accertare il corretto funzionamento delle unit\`a di codice anche con input limite.
\item Prova del prodotto: sar\`a effettuata eseguendo l'applicazione e interagendo con essa, assicurandosi che corrisponda a quanto progettato.
\item Confronto con altri software gi\`a esistenti: requisito fondamentale del prodotto \`e la capacit\`a di generare immagini in formato \underline{SVG}. Per assicurarsi che i file generati siano corretti, si caricheranno su software la cui correttezza \`e nota (ad esempio \textit{Inkscape}) e ci si assicurer\`a che le immagini visualizzate corrispondano con quanto disegnato dall'utente. 
\item Lettura del codice per controllare che le norme di codifica siano rispettate
\item Analisi del codice tramite \textit{walkthrough}, alla ricerca di eventuali errori. Vista la particolare dispendiosit\`a di questa modalit\`a di verifica, essa verr\`a applicata unicamente a sezioni critiche del codice.
\end{elenconumerato}


\subsezione{Risorse necessarie, risorse disponibili}
\paragraph{Documenti} Ogni documento sar\`a verificato da un revisore. Sarebbe auspicabile che a documenti particolarmente importanti, come ad esempio la definizione di prodotto, venissero assegnati due revisori, ma i costi eccessivi associati a una tale scelta non lo permettono.
\paragraph{Codice} Il numero di verificatori da assegnare ad un dato modulo di codice varieranno in base alle esigenze del modulo stesso.

\subsezione{Strumenti, tecniche, metodi}
Per il tracciamento degli errori e delle loro soluzioni ci si avvarr\`a degli strumenti di bug tracking messi a disposizione da \textit{SourceForge} e dagli strumenti di versionamento in uso.
Le discussioni tra i verificatori e gli autori delle modifiche avverrano tramite il sistema di commenti di \textit{GitHub} qualora si tratti di annotazioni di scarsa entit\`a, mentre discussioni pi\`u rilevanti avverranno nel gruppo di discussione interno all'azienda. Entrambi i sistemi mantengono in modo persistente le informazioni, permettendo dei riferimenti futuri alle stesse. I verificatori saranno avvisati della necessit\`a di un'attivit\`a di verifica tramite un messaggio nel gruppo di discussione interno. 


\sezione{Gestione amministrativa della revisione}

\subsezione{Comunicazione e risoluzione di anomalie}
Quando un'anomalia viene riscontrata durante una verifica, il verificatore provveder\`a ad aggiungere un ticket al sistema di bug tracking e a segnalarlo, usando i metodi sopra esposti, a chi dovr\`a correggere il problema. La segnalazione dell'errore dovr\`a contenere una descrizione chiara del problema riscontrato. Qualora il correttore trovasse poco chiara o ambigua la descrizione, o fosse in disaccordo con la stessa, aprir\`a una discussione con il verificatore nel gruppo di discussione interno. A correzione avvenuta, questo sar\`a segnalato al verificatore tramite il gruppo di discussione. Il verificatore applicher\`a nuovamente il procedimento di verifica alla correzione. Se quest'ultimo passo va a buon fine, il correttore porr\`a come risolto il ticket nel sistema di bug tracking.

\subsezione{Trattamento delle discrepanze}
Qualora un prodotto si discostasse da quanto descritto nell'analisi dei requisiti, il responsabile della qualit\`a provveder\`a a individuare la causa e la persona, o le persone, imputabili per la discrepanza. Il responsabile della qualit\`a, dopo averne discusso con lui, o loro, traccier\`a delle azioni correttive e si assicurer\`a che vengano attuate.

\subsezione{Procedure di controllo di qualit\`a di processo}
L'amministratore del progetto deve assicurarsi che gli standard descritti nelle \textit{NormeDiProgetto} siano rispettati, e deve segnalare eventuali infrazioni al responsabile, in modo che provveda a risolverle. Il responsabile della qualit\`a controlla che le attivit\`a di verifica vengano svolte, e che seguano le procedure descritte in questo documento.
\newpage

\sezione{Resoconto delle attivit\`a di verifica}

\subsezione{Tracciamento Aspettative - Requisiti}
\begin{table}[h]
\begin{center}
     \begin{tabular}
           {@{\extracolsep{\fill}}|c|c|}
     \hline
%%%%%%%%%%%%%%INTESTAZIONE COLONNE%%%%%%%%%%%%%%%%%%%%%%%%%%%%%%%%
      \textbf{Aspettative} & \textbf{Requisiti} \\
%%%%%%%%%%%%%%FINE INTESTAZIONE COLONNE%%%%%%%%%%%%%%%%%%%%%%%%%%%%%%%%%%%%%
      \hline
     ACO-1 & RPO-1 RQO-9 \\
     \hline
     ACO-2 & RFO-1 RFO-2 RFO-3 RFO-4\\ & RFO-5 RFO-6 RFO-7 RFO-8\\ & RFO-11 RF0-12 RFO-13 RFO-14 \\
     \hline
     ACO-3 & RFO-8 RFO-9 RFO-10 \\
     \hline
     ACO-4 & RFO-15 RFO-16 \\
     \hline
     ACO-5 & RQO-1 RQO-7 RQO-8 RAO-1 RAO-2 RAO-6 \\
     \hline
     ACO-6 & RAO-3 \\
     \hline
     ACO-7 & RIO-1 RIO-2 RIO-3 RIO-4 \\
     \hline
     ACO-8 & RAO-4 RQO-2 RQO-3 RQO-4 RQO-5 RQO-6 \\
     \hline
     ACD-1 & RAO-5 \\
     \hline
     ACF-1 & RFD-3 RFD-4 \\
     \hline
     ACF-2 & RFD-1 RFD-2 RFF-1 RFF-2 \\
     \hline
     ACF-3 & RAD-1 \\
%%%%%%%%%%% PARTE DA MODIFICARE %%%%%%%%%%%%%%%%
    \hline %%FINE RIGA
%%%%%%%%%%% FINE PARTE DA MODIFICARE %%%%%%%%%%%%%%%%%%%%%%%%%%%%%%%%%%%%%%%%
    \end{tabular}
  \caption{Aspettative - Requisiti} %INSERIRE DIDASCALIA - SE NECESSARIA -
  \label{tab:aspettative}
  \end{center}
\end{table}
 \newpage
 
 \subsezione{Tracciamento Use Case - Requisiti}
\begin{table}[h]
\begin{center}
     \begin{tabular}
           {@{\extracolsep{\fill}}|c|c|}
     \hline
%%%%%%%%%%%%%%INTESTAZIONE COLONNE%%%%%%%%%%%%%%%%%%%%%%%%%%%%%%%%
      \textbf{Aspettative} & \textbf{Requisiti} \\
%%%%%%%%%%%%%%FINE INTESTAZIONE COLONNE%%%%%%%%%%%%%%%%%%%%%%%%%%%%%%%%%%%%%
      \hline
     UC0 & RA0-1 \\
     \hline
     UC1 & RFD-3  \\
     \hline
     UC2 & RIO-2  RIO-4\\
     \hline
     UC3 & RIO-3 \\
     \hline
     UC4 & RFO-13 \\
     \hline
     UC5 UC6 & RFO-14 \\
      \hline
     UC7 UC8 & RFO-15 RFO-16 \\
     \hline
     UC9.1 & RFO-1 RFO-2 RFO-9\\
     \hline
     UC9.2 &  RFO-3 RFO-9\\
     \hline
     UC9.3 & RFO-4 RFO-9\\
     \hline
     UC9.4 & RFO-5 RFO-9\\
     \hline
     UC9.5 & RFO-6 RFO-7 RFO-9\\
     \hline
     UC9.6 UC9.7 & RFO-9 \\
     \hline
     UC9.9 & RFF-1 RFO-9 \\
     \hline
     UC9.10 & RFD-1 RFO-9\\
     \hline
     UC9.11 & RFF-2 RFO-9 \\
     \hline 
     UC9.12 & RFD-2 RFO-9 \\
     \hline
     UC 9.8 UC10.1 UC10.2 &  RFO-7 \\
     \hline
     UC10.3 UC10.4 UC10.5 & RFO-9 \\
     \hline
     UC10.5 UC10.7 UC10.10 &  RFO-8 \\
     \hline
     UC10.8 UC10.9 & RFO-10 \\
     \hline
     UC10.11 & RFD-4 \\
     
%%%%%%%%%%% PARTE DA MODIFICARE %%%%%%%%%%%%%%%%
    \hline %%FINE RIGA
%%%%%%%%%%% FINE PARTE DA MODIFICARE %%%%%%%%%%%%%%%%%%%%%%%%%%%%%%%%%%%%%%%%
    \end{tabular}
  \caption{Use Case - Requisiti} %INSERIRE DIDASCALIA - SE NECESSARIA -
  \label{tab:aspettative}
  \end{center}
\end{table}
 \newpage
 
 
\subsezione{Dettaglio delle verifiche tramite analisi}
\subsezione{Dettaglio delle verifiche tramite prove (test)}
\subsezione{Dettaglio dell'esito delle revisioni}


\sezione{Pianificazione ed esecuzione del collaudo}

\subsezione{Specifica della campagna di validazione (collaudo incluso)}
La campagna di validazione si svolger\`a in due fasi distinte:
\begin{elenconumerato}[\textbf{}]{\subsubsecindent}
\item \textit{Alpha test}: verr\`a eseguito internamente all'azienda. Il test consister\`a nell'esecuzione dell'applicativo per controllare che quanto descritto nell'analisi dei requisiti, la quale ha valore contrattuale, sia presente nel prodotto. Il test sar\`a guidato dagli \textit{use case} presenti nell'analisi dei requisiti.
\item \textit{Beta test}: il prodotto verr\`a distribuito al committente a fini di test.
\end{elenconumerato}
\subsezione{Dettaglio dell'esito della campagna di validazione}

\end{document}
